\section{下一步：lossy systems、异构硬件与 agentic engineering}

课程最后把 correctness 从模型供应商重新交回应用。共享 provider 面向许多未知 workload，必须保守；单应用部署知道自己的 eval、prompt 与错误成本，可以尝试更激进的 layer skip、pruning、heavy quantization 或 lossy speculation。与此同时，draft model 与特殊 layer 可能被编成 megakernel，prefill/decode 也可能迁移到不同加速器。

\subsection{应用定义 correctness，才有资格做 lossy optimization}

所谓 “unhinged” optimization 不是忽略质量，而是用应用自己的 eval 明确允许损失。若某类 data extraction 能容忍极小字段错误并由规则校验，就可能接受更激进 quantization；若是医疗、金融或安全 tool call，质量预算可能接近零。

\lecturefigure{slide-069.jpg}{未来的 lossy optimization、megakernel 与新硬件}{官方 deck 第 69 页；课堂 01:19:56--01:20:44}

\begin{knowledgebox}{读图：三条未来路线共享“specialization”}
Lossy optimization 针对单应用正确性；megakernel 针对固定模型层或 draft workload；异构硬件针对 prefill/decode 不同算术强度。它们都用更窄适用范围换取更高效率，因此必须同时投资 compiler、routing、fallback 与 eval。
\end{knowledgebox}

\begin{warningbox}{预测不是承诺}
讲者关于 prefill/decode 使用不同 accelerator、analog/optical/memristive computing 的判断属于未来方向。硬件可用性、软件栈、精度、网络与经济性都未由这张 slide 证明，不能写进当前容量计划的 guaranteed roadmap。
\end{warningbox}

\subsection{Agent 不会取消 correctness engineering}

如果 agent 负责 benchmark、configuration、on-call 或写 bespoke engine，工程师的工作会转向定义正确性、权限、工具和反馈。把 AWS root account 交给 agent 并不会自动产生可靠 infra；必须提供可验证目标、sandbox、可回滚操作、metrics、trace 与 human escalation。

\lecturefigure{slide-071.jpg}{Agentic software engineering 的核心仍是 correctness system}{官方 deck 第 71 页；课堂 01:20:50--01:22:19}

\begin{knowledgebox}{读图：现在与以后分别需要什么}
近期 agent 可以辅助 benchmark、配置搜索与监控，但前提是系统有清晰指标和权限边界；更长期的设想是按模型--硬件--workload 自动生成 bespoke engine。VibeServe（arXiv:2605.06068）在课堂日前展示了多 agent 搜索、实现、correctness check 与性能测量的闭环，但它仍是研究证据，不是无需监督的生产替代品。
\end{knowledgebox}

\begin{lstlisting}[caption={Agentic inference engineering 的安全闭环}]
while budget_remaining:
    candidate = agent.propose_change(
        workload, source_code, allowed_tools, constraints
    )
    build_in_sandbox(candidate)
    if not correctness_suite_passes(candidate):
        reject(candidate)
        continue
    metrics = benchmark(candidate, production_trace_sample)
    if dominates_baseline(metrics) and rollback_is_ready():
        canary(candidate)
    else:
        reject(candidate)
\end{lstlisting}

\begin{importantbox}{CI/CL 的合理解释}
Deck page 072 提出 continuous integration / continuous learning loop，但实际录像在 page 071 后结束，因此本讲义不把该页当成课堂已讲内容。作为延伸，可把它理解为：更多用户产生更多经治理的数据，数据支持 eval、distillation/speculator training 与模型更新，改进 UX 后再产生新反馈。任何自动学习都必须有数据许可、质量 gate 与 rollback。
\end{importantbox}

\subsection{本章小结}

未来 serving 会更加 specialized，但 specialization 不等于放弃系统纪律。应用越了解自身正确性，越能安全使用 lossy optimization；agent 越能修改系统，越需要强 eval、observability、sandbox 和权限模型。自动化提高的是搜索与执行速度，工程责任转移到目标与验证。

